% !PRINTED 619
\weberpage{619}
%% [page furniture, not body text -- source-confirmed 2026-08-10]
%% \S\,169. \quad Beziehung der Teilungsk\"orper zu dem Klassenk\"orper. \quad 619

%<<CONT
und der Relativgrad in bezug auf $\Omega$ h\"ochstens gleich $\tfrac{1}{2} h''\varphi(m)$, und daraus geht nach (11) hervor, da\ss\ $h''$ nicht kleiner als $H'_g$ sein kann. Damit ist bewiesen:

5. \gesperrt{Die Klassengleichung $H_{-D}(u)$ ist in dem K\"orper, der alle Einheitswurzeln enth\"alt, nicht weiter zerlegbar als in die den Geschlechtern entsprechenden Faktoren.}

\noindent\textbf{\S\,169.\qquad \gesperrt{Beziehung der Teilungsk\"orper zu dem Klassenk\"orper.}}

Wir haben in den beiden letzten Paragraphen zwei Arten von Gruppen betrachtet, von denen die erste zu den Teilungsk\"orpern der elliptischen Funktionen, die zweite zu den Ordnungsk\"orpern in Verbindung mit Einheitswurzeln f\"uhrte. Zwischen diesen besteht eine Beziehung:

Die erste Gruppe $A$ bestand aus den Zahlen $\alpha$ in $O$, die der Bedingung
\weq{1}{\alpha \equiv \pm 1 \pmod{\mathfrak{m}}}
gen\"ugten, wenn $\mathfrak{m}$ irgend ein Ideal in $\Omega$ ist, die zweite Gruppe $A'$ bestand aus den Zahlen $\alpha'$ der Ordnung $[Q]$, die der Bedingung
\weq{2}{N(\alpha') \equiv 1 \pmod{m}}
gen\"ugten, wenn $m$ und $Q$ nat\"urliche Zahlen bedeuten.

Nehmen wir an, es sei $\mathfrak{m}$ ein Primideal oder eine Potenz eines Primideals, so k\"onnen wir $Q$ und $m$ so annehmen, da\ss\ die Gruppe $A'$ in der Gruppe $A$ enthalten ist.

Denn ist $r$ eine rationale Zahl, so ist, da $\alpha'$ der Ordnung $[Q]$ angeh\"ort:
\weq{3}{\alpha' \equiv r \pmod{Q}, \quad N(\alpha') \equiv r^2 \pmod{Q},}
und wenn $Q$ durch $m$ teilbar ist:
\weq{4}{\alpha' \equiv r \pmod{m},}
also nach (2) und (3)
\[
r^2 \equiv 1 \pmod{m}
\]
\weq{5}{(r - 1)(r + 1) \equiv 0 \pmod{m},}
wir nehmen also $Q$ durch $m$ und $m$ durch $\mathfrak{m}$ teilbar an; dann ist nach (5), wenn $\mathfrak{m}$, wie wir angenommen haben, ein Primideal oder eine Potenz eines solchen ist, entweder
\weq{6}{r \equiv 1 \text{ oder } r \equiv -1 \pmod{\mathfrak{m}},}
%CONT>>
